\expandafter\def\csname HistoricalNoteText18\endcsname{Since this value differs, from the twentieth digit onward, from that calculated by the distinguished \textsc{Mascheroni} in Adnotat. ad \textsc{Euleri} \textit{Calculum Integr.}, I urged \textsc{Fridericus Bernhardus Gothofredus Nicolai}, a young man tireless in calculation, to repeat that calculation and extend it further. He thus found, by two calculations, namely working down both from $z=50$ and from $z=100$,
\[
\Psi0=-0{,}5772156649\ 0153286060\ 6512090082\ 4024310421
\]
The second part of the table appended at the end of this section is also due to the same very practised calculator; it gives the values of the function $\Psi z$ to 18 digits (the last of which is uncertain), for all values of $z$ from 0 to 1 at steps of one hundredth. The methods by which both tables were constructed rest partly on the theorems given here, partly on special devices of calculation, which we shall present on another occasion.}
